% !TEX TS-program = lualatex

\documentclass{standalone}

%%%
% Des couleurs faciles d'emploi via le package \xcolor!
%%%
\usepackage[svgnames]{xcolor}

%%%
% La bibliothèque ''luadraw'' allie une facilité d’utilisation
% à un rendu particulièrement soigné.
%%%
\usepackage[3d]{luadraw}

\begin{document}

\begin{luadraw*}{name = angular-gradient-easy-way}
local ld = luadraw

local palext = require 'luadraw_palettes'

------
-- La bibliothèque ''luadraw_shadedforms'' proposée par ¨luadraw
-- permet d'obtenir des dégardés facilement.
--
-- note::
--     ''luadraw_shadedforms'' charge ''luadraw_palettes'' pour
--     nous.
------
require 'luadraw_shadedforms'

local graphview = ld.graph:new{
  size = {10, 10},
  bbox = false
}

------
-- Le principe repose sur la ¨def d'un cercle et d'une ¨fonc de
-- gradient associée. Le choix d'un tracé centré en `0` et d'une
-- variation de couleur indexée sur les \arg d'affixes complexes
-- permet d'obtenir le rendu annulaire souhaité.
--
-- note::
--     Nous n'utlisons pas ''C = circle(0, 2)'', car ceci crée
--     un effet de palier non paramétrable à la main.
------
C = ld.polar(
  function(t)
    return 2
  end,
  -math.pi, math.pi,
  200
)

local f = function(x,y)
  return ld.cpx.arg(ld.cpx.Z(x,y))
end

------
-- Déléguons le travail à ¨luadraw.
------
graphview:Dshadedpolyline(
  C,              -- La ligne polygonale.
  palext.Picnic,      -- La palette.
  {
    values = f,   -- Calcul du gradient.
    width  = 400  -- Épaisseur excessive.
  }
)

------
-- Montrons le résultat de notre oeuvre.
------
graphview:Show()
\end{luadraw*}

\end{document}
